\documentclass[11pt, letterpaper]{article}

% ---------- packages ----------
\usepackage[margin=0.7in]{geometry}
\usepackage{tcolorbox}
\usepackage{enumitem}
\usepackage{hyperref}
\usepackage{xcolor}
\usepackage{titlesec}
\usepackage{fancyhdr}
\usepackage{graphicx}
\usepackage{tabularx}
\usepackage{booktabs}

% ---------- colors ----------
\definecolor{boxborder}{RGB}{60, 90, 150}
\definecolor{boxbg}{RGB}{245, 248, 255}
\definecolor{headcolor}{HTML}{00636f}
\definecolor{placeholder}{RGB}{140, 140, 140}

% ---------- tcolorbox style ----------
\tcbset{
  sectionbox/.style={
    colframe=boxborder,
    colback=boxbg,
    boxrule=0.6pt,
    arc=2pt,
    left=6pt, right=6pt, top=4pt, bottom=4pt,
    fontupper=\small,
  }
}

% ---------- placeholder command ----------
\newcommand{\placeholder}[1]{\textcolor{placeholder}{\textit{#1}}}

% ---------- section formatting ----------
\titleformat{\section}
  {\normalsize\bfseries\color{headcolor}}{\thesection.}{0.4em}{}
\titlespacing*{\section}{0pt}{6pt}{2pt}

% ---------- header / footer ----------
\pagestyle{fancy}
\fancyhf{}
\renewcommand{\headrulewidth}{0.4pt}
\lhead{\small\color{headcolor} CS 4100 -- Project Executive Summary}
\cfoot{\small\thepage}

% ---------- suppress paragraph indent ----------
\setlength{\parindent}{0pt}
\setlength{\parskip}{3pt}

\begin{document}

% ===== TITLE BLOCK =====
\begin{center}
  {\Large\bfseries\color{headcolor} Project Title Here}\\[4pt]
  {\normalsize Member 1, Member 2, Member 3, Member 4}\\[2pt]
\end{center}

\vspace{4pt}

% ===== 1. MOTIVATION & PROBLEM STATEMENT =====
\section{Motivation \& Problem Statement}
\begin{tcolorbox}[sectionbox, height=4.6cm, valign=top]
  \placeholder{%
    Describe the real-world problem your project addresses and explain
    why it matters. Identify the specific gap, limitation, or
    opportunity that motivates this work. State the core research
    question or engineering goal in one clear sentence.
    Cite relevant prior work or context that frames the problem
    (1--2 references are sufficient). Aim for roughly 8--12 sentences.}
\end{tcolorbox}

% ===== 2. APPROACH =====
\section{Approach}
\begin{tcolorbox}[sectionbox, height=7.0cm, valign=top]
  \placeholder{%
    Summarize your technical approach at a high level. Describe the
    system architecture or methodology, key algorithms or models used,
    and the most important design decisions. If applicable, include a
    small figure or diagram (use \textbackslash includegraphics), and resize to stay within bounded box area.
    Mention the primary tools, frameworks, or libraries
    (e.g., PyTorch, Mesa, Gymnasium, PettingZoo).
    Clearly state any assumptions or scope limitations.
    Aim for roughly 10--15 sentences or equivalent with an optional figure.}
\end{tcolorbox}

% ===== 3. KEY RESULTS =====
\section{Key Results}
\begin{tcolorbox}[sectionbox, height=8cm, valign=top]
VGG16 achieved 99.63\% validation accuracy and 99.55\% ROC-AUC on StyleGAN1 using pretrained VGGFace weights, fine-tuned over 15 epochs at a learning rate of 0.0001. ResNet also performed strongly in the within-GAN setting. Noise pattern analysis hit 91.72\% on StyleGAN1 using an XGBoost classifier on 12,297 noise residual features per image; when tested on StyleGAN2 it dropped to 62\%, and on diffusion-generated images (CIFAKE and DeepFakeFace) collapsed to 46--50\%, essentially random. The noise fingerprint a model learns is specific to the generator it trained on.

DCT with gradient boosting showed better cross-GAN behavior than expected; while within-GAN accuracy was lower than VGG and noise pattern, it held up more consistently on unseen generators, suggesting frequency-domain features are more generalizable than raw noise residuals. DIRE failed across all four datasets at 50--55\%. The cause was model-data mismatch; our DDPM was pretrained on 32$\times$32 CIFAR-10 images, not faces, so real and fake faces reconstructed equally poorly with no signal for the classifier. This does not invalidate DIRE; it works well with a domain-matched model, but is highly sensitive to that choice.

The overarching finding is consistent; within-GAN accuracy is high, cross-GAN accuracy drops significantly, and performance on diffusion models is near random for methods trained on GAN data. Deep learning models like VGG achieved the highest peak accuracy but the sharpest out-of-distribution drops. Methods that sacrificed peak accuracy, like DCT, actually generalized better. A reliable detector would need training on a large and diverse dataset spanning many GAN architectures and diffusion models simultaneously; without that breadth, any detector will have blind spots against generators it has never seen.
\end{tcolorbox}

\newpage

\begin{tcolorbox}[sectionbox, height=8cm, valign=top]
\includegraphics[width=0.42\linewidth]{Picture.png}\hfill
\includegraphics[width=0.42\linewidth]{Picture2.png}
  \\
  \placeholder{%
  Replace these figures with your own, or remove them and use this space for more text. Ensure all text in your figures are clear and readable at this scale, similar to these examples.
   }
\end{tcolorbox}

% ===== 4. CHALLENGES & LESSONS LEARNED =====
\section{Challenges \& Lessons Learned}
\begin{tcolorbox}[sectionbox, height=7.2cm, valign=top]
The most significant technical challenge was poor cross-GAN generalization across nearly every detection method. Models trained on StyleGAN2 achieved near-perfect accuracy on that dataset but dropped to around 50\% when tested on StyleGAN1 or Stable Diffusion, showing that models learned artifacts specific to one GAN rather than general signs of a fake image. The DCT + Gradient Boosting results made it clear that frequency-domain artifacts differ considerably depending on which GAN generated the image. Fine-tuning VGG was slow and demanding even on Kaggle GPUs, and completing 15 epochs took a considerable amount of time. The Vision Transformer was even harder to work with and the team was only able to complete one epoch of training; because one epoch was not enough for the model to learn meaningfully, the team decided not to include those results in the final comparison. This made it clear that transformer-based models require far more compute than CNNs and should only be attempted when stronger hardware is available. Coordinating experiments across four team members was also a challenge and required careful version control to keep model checkpoints and dataset splits consistent. A major lesson learned is that teams should agree on standardized evaluation splits and metrics early so that results from different methods can be compared fairly. Overall the project showed that a simpler, well-trained model will outperform a complex model that did not have enough time or resources to train properly.
\end{tcolorbox}

% ===== 5. INDIVIDUAL CONTRIBUTIONS =====
\section{Individual Contributions}
\begin{tcolorbox}[sectionbox, height=4.3cm, valign=top]
  \placeholder{Replace the table below with your team's actual contributions.}

  \vspace{4pt}
  \centering
  \begin{tabularx}{0.95\textwidth}{l X c}
    \toprule
    \textbf{Member} & \textbf{Primary Contributions} & \textbf{Approx.\ \%} \\
    \midrule
    Member 1 & \placeholder{e.g., System design, agent logic}          & 20\% \\
    Member 2 & \placeholder{e.g., Data pipeline, evaluation scripts}   & 20\% \\
    Member 3 & \placeholder{e.g., Environment design and implementation} & 20\% \\
    Member 4 & \placeholder{e.g., Data pipeline, evaluation scripts}  & 20\% \\
    Member 5 & \placeholder{e.g., System design, agent logic}  & 20\% \\
    \bottomrule
  \end{tabularx}
\end{tcolorbox}

% ===== 6. REPOSITORY & REFERENCES =====
\section{Repository \& References}
\begin{tcolorbox}[sectionbox, height=4.5cm, valign=top]
  \textbf{GitHub Repository:}
  \placeholder{\href{https://github.com/your-org/your-repo}%
       {\texttt{https://github.com/your-org/your-repo}}}\\
    \placeholder{Ensure repo is public, or add instructor as a collaborator.}\vspace{1pt}

  \textbf{Project Tracker:} \placeholder{\href{https://github.com/your-org/your-repo}{\texttt{https://github.com/your-org/projects/your-project}}}

  \vspace{6pt}
  \textbf{Key References:}
  \begin{enumerate}[leftmargin=1.4em, itemsep=1pt, topsep=2pt, label={[\arabic*]}]
    \item \placeholder{Author(s), ``Title,'' Venue, Year. Brief annotation.}
    \item \placeholder{Author(s), ``Title,'' Venue, Year. Brief annotation.}
    \item \placeholder{Author(s), ``Title,'' Venue, Year. Brief annotation.}
  \end{enumerate}

  \vspace{4pt}
  \placeholder{Add or remove reference entries as needed (3--4 recommended).}
\end{tcolorbox}

\end{document}
